\section{Related Work}

Vibration-based fault diagnosis has evolved significantly, from traditional signal processing to advanced deep learning and multimodal approaches. This section reviews key developments, highlighting limitations in domain generalization and zero-shot capabilities that our VibSemAlign addresses.

\subsection{Vibration Fault Diagnosis}
\subsubsection{Traditional Signal Processing}
Traditional methods rely on handcrafted features from time, frequency, and time-frequency domains. Fast Fourier Transform (FFT) decomposes signals into frequency components, identifying fault-specific harmonics like ball pass frequencies (BPFO/BPFI) \cite{cite1}. Envelope analysis demodulates amplitude-modulated fault impulses, while wavelet transforms capture transient non-stationarities \cite{cite2}. These techniques are foundational in aerospace and automotive industries where interpretability is paramount. However, real-world complications like speed variations require order tracking or Vold-Kalman filtering as preprocessors.

Empirical Mode Decomposition (EMD) and Variational Mode Decomposition (VMD) adaptively decompose signals into intrinsic mode functions (IMFs) \cite{cite3}. VMD outperforms EMD by optimizing mode bandwidth variationally, reducing mode mixing. Both excel in interpretability but struggle with noise, variable speeds, and domain shifts \cite{cite10}. EMD suffers from mode mixing where disparate oscillations contaminate neighboring IMFs, and end effects introducing spurious components. VMD alleviates these via variational optimization yet requires careful tuning of mode number $K$ and balancing parameter $\alpha$. Under domain shifts---such as load variations altering fault frequencies---decompositions yield misaligned IMFs, diminishing feature robustness.

Complementary techniques include spectral kurtosis for impulsivity quantification, minimum entropy deconvolution (MED) for enhancing fault impulses, and cyclostationary analysis for modulated fault signatures. Spectral kurtosis identifies frequency bands with maximal impulsivity, enabling optimal bandpass filter design for envelope analysis. MED and its variant Maximum Correlated Kurtosis Deconvolution (MCKD) iteratively deconvolve the signal to emphasize periodic fault impulses buried in noise. Cyclostationary analysis exploits the periodic modulation structure inherent in rotating machinery, revealing fault signatures invisible to conventional FFT. These methods collectively provide a rich toolkit for signal characterization but share the fundamental limitation of handcrafted feature design: inability to adapt to distributional shifts without manual recalibration.

\subsubsection{Deep Learning Approaches}
Deep learning revolutionized fault diagnosis by automating feature learning. Convolutional Neural Networks (CNNs) process 1D time-series or 2D spectrograms, capturing local patterns like fault impulses and harmonics \cite{cite12,cite13}. Recurrent Neural Networks (RNNs/LSTMs) model temporal dependencies, while Transformers leverage self-attention for long-range interactions \cite{cite14}. Residual architectures with disentangled representation learning further improve feature extraction \cite{cite37}. End-to-end models like DCNNWK and Multi-Scale CNNs handle multi-fault scenarios \cite{cite35}. Graph Neural Networks (GNNs) represent signals as graphs for relational reasoning \cite{cite36}.

These achieve 95--99\% accuracy on source domains (e.g., CWRU) but degrade 20--40\% in cross-domain settings without retraining \cite{cite17,cite31}. CNNs overfit to source-domain impulse signatures and fail to capture shifted harmonic patterns; Transformers disperse attention across domain-specific artifacts. Data hunger remains critical: DL demands thousands of labeled samples per fault/domain, infeasible for rare industrial faults. Augmentation (e.g., GANs) helps but introduces artifacts that may not faithfully represent the underlying fault physics.

Recent works explore Vision Transformers (ViTs) for vibration spectrograms, treating time-frequency representations as patch sequences \cite{cite31}. ViTs capture global spectral dependencies through self-attention, outperforming CNNs on cross-load scenarios by 3--5\%. Vision Transformers combined with frequency-domain fusion achieve competitive results on standard benchmarks \cite{cite41}. However, they remain vulnerable to domain shift without explicit alignment mechanisms, as self-attention weights distribute across both diagnostically relevant and irrelevant spectrogram regions.

\subsection{Multimodal Learning and VLMs}
\subsubsection{Contrastive Alignment (CLIP-like)}
Contrastive Language-Image Pretraining (CLIP) aligns image/text embeddings via InfoNCE loss, enabling zero-shot classification through semantic similarity \cite{cite20}. CLIP's success spawned adaptations: ALIGN scales to billions of pairs; OpenCLIP refines architectures \cite{cite27,cite28}. Multi-modal contrastive learning frameworks have demonstrated promise in bridging signal-language representations \cite{cite39}. In diagnostics, spectrograms treated as ``images'' can align with fault descriptions, but vibrations' unique spectral traits demand modality-specific tuning absent in vanilla CLIP \cite{cite29}.

The fundamental challenge lies in the semantic gap between natural images and vibration spectrograms. Natural images exhibit spatial locality, rich textures, and object semantics that CLIP's visual encoder learns through extensive pretraining. Vibration spectrograms, conversely, display periodic harmonic structures, modulation sidebands, and transient impulses whose diagnostic significance depends on precise frequency localization---patterns fundamentally different from visual object recognition. This modality mismatch explains why vanilla CLIP achieves only 42\% zero-shot accuracy on vibration fault classification versus 76\% on ImageNet.

\subsubsection{VLMs for Non-RGB Modalities}
Vision-Language Models like LLaVA and BLIP-2 extend CLIP with generative decoders for visual question answering \cite{cite21}. Non-RGB applications include medical imaging (MRI/CT-text) and remote sensing (SAR-radar descriptions). Recent works integrate LLMs with domain-specific visual encoders: BiomedCLIP for medical imaging, RemoteCLIP for satellite imagery, demonstrating that modality-adapted VLMs substantially outperform generic counterparts in specialized domains.

In fault diagnosis, MMFault and LLM4FD feed spectrograms to VLMs with diagnostic prompts \cite{cite23,cite8}. VLMIFD and MMIFD incorporate industrial priors; FaultGPT generates synthetic labels \cite{cite24,cite25,cite26}. Large language models (LLMs) have emerged as complementary tools for maintenance knowledge representation and fault reasoning \cite{cite40}. Recent 2024--2025 works explore multi-agent LLM frameworks for maintenance planning, retrieval-augmented generation (RAG) for fault knowledge bases, and chain-of-thought reasoning for root cause analysis. Yet existing approaches overlook dedicated vibration-text alignment, relying on image-biased VLMs suboptimal for spectrograms. MMFault falters on noisy industrial spectrograms by mistaking modulation sidebands for image textures; VLM-IFD demonstrates prompt sensitivity and struggles with compound faults. None achieves robust zero-shot cross-domain performance without target-domain adaptation.

\subsection{Domain Adaptation and Zero-Shot Learning}
\subsubsection{Transfer Learning in Diagnostics}
Domain adaptation aligns source/target distributions via adversarial training (DANN), Maximum Mean Discrepancy (MMD), or CORAL \cite{cite34,cite22,cite10}. In vibrations, transfer from CWRU to Paderborn yields 10--15\% gains but requires target (unlabeled) data \cite{cite17}. Cycle-consistent generative adversarial networks (CycleGAN) enable cross-domain signal translation \cite{cite38}. Recent conditional domain adaptation methods address class-conditional shifts, while partial transfer learning handles asymmetric label spaces between domains. However, all require access to target domain data during training, limiting applicability to truly unseen operational environments.

Zero-shot learning (ZSL) for diagnostics uses auxiliary attributes (fault severity, component type) to bridge seen/unseen classes \cite{cite33}. Attribute-based ZSL maps signals to semantic descriptors, enabling classification of novel faults. However, attribute design is labor-intensive and domain-specific, limiting scalability. Generalized ZSL (GZSL) addresses the more realistic scenario where both seen and unseen classes appear at test time, but performance degrades substantially (10--20\%) compared to standard ZSL.

\subsubsection{Few-Shot and Meta-Learning}
Few-shot learning adapts with 1--5 target shots, achieving 82--88\% accuracy on vibration fault classification \cite{cite31}. MAML optimizes for fast adaptation: $\min_\theta \mathbb{E} [\mathcal{L}_{task}(\theta'(\theta))]$ \cite{cite35}. Prototypical networks compute class prototypes in embedding space for nearest-neighbor classification. Relation networks learn similarity metrics end-to-end. These show promise for vibration applications but falter without semantic guidance---embedding spaces learned from signal features alone lack the rich semantic structure that language provides.

\subsection{Research Gap and Our Contribution}

Existing approaches exhibit three critical gaps: (1) no dedicated vibration-semantic alignment mechanism---vanilla CLIP's image-text pretraining poorly transfers to spectrogram modality; (2) zero-shot cross-domain performance remains below 85\% even with domain adaptation techniques requiring target data; (3) few-shot methods lack semantic grounding, limiting transferability across machinery types. VibSemAlign bridges these gaps through spectrogram-specific contrastive fine-tuning, physics-informed prompt engineering, and MAML-based fast adaptation, achieving 96.5\% zero-shot CWRU accuracy and 85.3\% cross-domain CWRU$\to$Paderborn without target labels---establishing a new paradigm for semantic-driven industrial diagnostics.
